\chapter{Зачем нужны контейнеры}

\section{Проблема, которую решают контейнеры}

Почти каждый разработчик хотя бы раз слышал фразу \enquote{у меня работает}.
Программа успешно запускается на ноутбуке автора и падает на сервере
или у коллеги. Причины обычно прозаичны: другая версия интерпретатора,
не установлена системная библиотека, отличается локаль, не хватает прав
на каталог, в переменных окружения остался след от старого проекта.

Список зависимостей приложения давно вышел за пределы файла с пакетами.
Ему нужны:

\begin{itemize}
  \item конкретные версии языка и библиотек;
  \item системные пакеты (например, \file{libpq} для PostgreSQL
        или \file{libmagic} для определения типов файлов);
  \item утилиты вроде \cmd{curl}, \cmd{ffmpeg}, \cmd{fonts-dejavu};
  \item настройки операционной системы: локаль, часовой пояс, лимиты
        на число открытых файлов;
  \item конфигурация и секреты, которые не должны попадать в репозиторий.
\end{itemize}

Чем больше компонентов, тем сильнее среда \enquote{дрейфует}: сервер,
настроенный год назад вручную, уже не совпадает с документацией,
а воспроизвести его в точности не может никто. Классические инструменты
управления конфигурацией (Ansible, Puppet, Chef) частично решают проблему,
но они описывают \emph{как} настроить машину, а не \emph{что} в итоге
должно работать. Результат зависит от исходного состояния сервера
и порядка применения изменений.

Контейнеры предлагают другой подход: приложение поставляется вместе
со своей средой в виде единого артефакта~--- образа. На любой машине
с Docker этот образ запускается одинаково, потому что внутрь уже упаковано
всё необходимое, кроме ядра операционной системы.

\section{Виртуальные машины: честная, но тяжёлая изоляция}

Исторически первую попытку решить проблему сделали виртуальные машины.
Гипервизор (KVM, Xen, VMware, Hyper-V) эмулирует компьютер целиком:
виртуальная машина получает собственное ядро, драйверы и полноценную
операционную систему.

\begin{figure}[htbp]
\centering
\begin{tikzpicture}[
  font=\sffamily\scriptsize,
  every node/.style={align=center},
  base/.style={draw=accent!45, fill=accent!7, rounded corners=1.5pt,
               minimum width=3.3cm, minimum height=0.62cm},
  guest/.style={draw=accent2!60!black, fill=accent2!10, rounded corners=1.5pt,
               minimum width=1.5cm, minimum height=0.62cm},
  engine/.style={draw=warm!65!black, fill=warm!10, rounded corners=1.5pt,
               minimum width=3.3cm, minimum height=0.62cm},
]
  % --- Виртуальные машины ---
  \node[font=\sffamily\bfseries\small, text=accent] at (2.0,5.25) {Виртуальные машины};
  \node[base] (hw1) at (2.0,0.4) {аппаратное обеспечение};
  \node[base, above=1.5pt of hw1] (os1) {хостовая ОС};
  \node[base, above=1.5pt of os1] (hyp) {гипервизор};
  \node[guest, above=1.5pt of hyp, xshift=-0.87cm] (v1) {гостевая ОС\\+ приложение};
  \node[guest, above=1.5pt of hyp, xshift=0.87cm] (v2) {гостевая ОС\\+ приложение};
  \draw[decorate, decoration={brace, amplitude=3pt}, draw=muted]
    ($(v1.north west)+(-0.05,0.08)$) -- ($(v2.north east)+(0.05,0.08)$)
    node[midway, above=3pt, text=muted] {полноценные ОС};

  % --- Контейнеры ---
  \node[font=\sffamily\bfseries\small, text=accent] at (8.2,5.25) {Контейнеры};
  \node[base] (hw2) at (8.2,0.4) {аппаратное обеспечение};
  \node[base, above=1.5pt of hw2] (os2) {хостовая ОС, общее ядро};
  \node[engine, above=1.5pt of os2] (eng) {Docker Engine};
  \node[guest, above=1.5pt of eng, xshift=-0.87cm] (c1) {приложение\\+ библиотеки};
  \node[guest, above=1.5pt of eng, xshift=0.87cm] (c2) {приложение\\+ библиотеки};
  \draw[decorate, decoration={brace, amplitude=3pt}, draw=muted]
    ($(c1.north west)+(-0.05,0.08)$) -- ($(c2.north east)+(0.05,0.08)$)
    node[midway, above=3pt, text=muted] {обычные процессы};
\end{tikzpicture}
\caption{Виртуальные машины несут собственное ядро и полную ОС,
а контейнеры используют ядро хоста и содержат только приложение
с библиотеками}\label{fig:vm-vs-ct}
\end{figure}

Плюсы виртуальных машин очевидны: сильная изоляция (своё ядро),
возможность запускать разные операционные системы, привычные инструменты
резервного копирования и миграции. Минусы не менее очевидны:

\begin{itemize}
  \item каждая ВМ занимает гигабайты диска и сотни мегабайт памяти;
  \item загрузка занимает десятки секунд;
  \item обновлять и патчить нужно каждую гостевую ОС;
  \item плотность размещения ниже: на одном сервере помещается
        заметно меньше приложений.
\end{itemize}

\section{Идея контейнера}

Контейнер~--- это обычный процесс операционной системы, которому
средствами ядра создали иллюзию отдельного компьютера. Процесс видит
свою файловую систему, свои сетевые интерфейсы, свой список процессов
и свои лимиты ресурсов, хотя физически всё это принадлежит хосту.

\begin{defbox}[title={Определение: контейнер}]
\textbf{Контейнер}~--- изолированный набор процессов, запущенный из образа
на общем ядре операционной системы. Изоляцию обеспечивают пространства
имён (namespaces), ограничения ресурсов~--- контрольные группы (cgroups),
а отдельное представление файловой системы~--- механизм объединения слоёв
(overlayfs).
\end{defbox}

\begin{defbox}[title={Определение: образ}]
\textbf{Образ} (\emph{image})~--- неизменяемый шаблон файловой системы
и метаданных, из которого создаются контейнеры. Образ состоит из слоёв
и идентифицируется тегом (\cmd{nginx:1.27}) или хешем (\cmd{sha256:...}).
\end{defbox}

Слово \enquote{контейнер} подчёркивает метафору: приложение \enquote{упаковано}
вместе с зависимостями, как груз в стандартном морском контейнере.
Стандартизация важнее содержимого: любой кран, любой порт и любой поезд
умеют работать с контейнером одного стандарта, не зная, что внутри.
В мире программ эту роль играет формат OCI (Open Container Initiative):
один и тот же образ запускается Docker, Podman, containerd и Kubernetes.

\section{Сравнение виртуальных машин и контейнеров}

\begin{table}[htbp]
\centering
\small
\begin{tabularx}{\textwidth}{@{}lXX@{}}
\toprule
 & \textbf{Виртуальная машина} & \textbf{Контейнер} \\
\midrule
Ядро ОС & собственное у каждой ВМ & общее, ядро хоста \\
Изоляция & аппаратная, сильная & на уровне ядра (namespaces, cgroups) \\
Размер & гигабайты & десятки--сотни мегабайт \\
Запуск & десятки секунд & доли секунды--несколько секунд \\
Плотность & единицы--десятки на сервер & десятки--сотни на сервер \\
Гостевая ОС & нужна и требует обслуживания & не нужна, только библиотеки \\
Совместимость & любая ОС внутри & ядро и архитектура хоста \\
Основной сценарий & разные ОС, жёсткая изоляция, legacy & микросервисы, CI/CD, тиражирование среды \\
\bottomrule
\end{tabularx}
\caption{Ключевые различия виртуальных машин и контейнеров}\label{tab:vm-vs-ct}
\end{table}

Важно понимать: контейнеры не \enquote{лучше} виртуальных машин вообще.
Это разные инструменты. Часто их комбинируют: в облаке контейнеры
запускают внутри виртуальных машин, чтобы получить и плотность,
и сильную границу изоляции.

\begin{warnbox}
Контейнер не является границей безопасности того же уровня, что виртуальная
машина. Все контейнеры на хосте используют одно ядро: уязвимость ядра
или ошибка в настройке прав может затронуть соседей. Для недоверенного
кода используйте дополнительную изоляцию~--- отдельные ВМ, gVisor, Kata
Containers или rootless-режим.
\end{warnbox}

\section{Что даёт контейнеризация}

\begin{enumerate}
  \item \textbf{Воспроизводимость.} Один и тот же образ запускается
        одинаково на ноутбуке, в CI и в production. Меньше расхождений~---
        меньше \enquote{работает у меня}.
  \item \textbf{Быстрый старт и масштабирование.} Контейнер поднимается
        за секунды. Если нагрузка выросла, достаточно запустить ещё
        несколько экземпляров того же образа.
  \item \textbf{Плотность.} На одном сервере умещается значительно больше
        приложений, чем при использовании ВМ,~--- меньше затрат
        на инфраструктуру.
  \item \textbf{Единый артефакт поставки.} Образ собирается один раз
        и проходит через все среды: разработка, тестирование, staging,
        production. Меняются только конфигурация и секреты.
  \item \textbf{Изоляция зависимостей.} Два сервиса могут требовать разные
        версии одной библиотеки и не конфликтовать.
  \item \textbf{Удобство CI/CD.} Сборка, тесты и публикация выполняются
        в контейнере с известным окружением, а не на разделяемом
        агенте сборки.
  \item \textbf{Неизменяемая инфраструктура.} Вместо правки работающего
        сервера вы собираете новый образ и заменяете контейнер.
        Откат~--- это запуск предыдущего тега.
\end{enumerate}

\section{Границы применимости}

Контейнеры подходят не для всего. Полезно заранее знать ограничения.

\begin{itemize}
  \item \textbf{Графические приложения.} Запустить GUI-программу в контейнере
        можно, но проброс X11/Wayland и звука усложняет конфигурацию;
        для десктопных задач контейнеры обычно не нужны.
  \item \textbf{Состояние и данные.} Контейнер эфемерен: удалили контейнер~---
        пропал его записываемый слой. Базы данных и файлы пользователей
        требуют томов (глава~7) и осознанного резервного копирования.
        Контейнер не отменяет необходимость бэкапов.
  \item \textbf{Ядро и архитектура.} Контейнер Linux не запустится на ядре
        Windows напрямую, а образ для \cmd{amd64} не пойдёт на \cmd{arm64}
        без эмуляции. Мультиархитектурные образы решают вторую проблему
        (раздел~10.5).
  \item \textbf{Тяжёлые stateful-нагрузки.} Кластеры баз данных в контейнерах
        эксплуатируют, но это требует опыта: производительность диска,
        кворум, резервное копирование, обновления без простоя.
  \item \textbf{Полная изоляция недоверенного кода.} Как уже сказано,
        общее ядро~--- это компромисс между плотностью и безопасностью.
\end{itemize}

\section{Docker и его окружение}

\begin{description}[style=nextline]
  \item[Docker] Набор инструментов и формат упаковки, который сделал
        контейнеры массовыми. В книге под Docker понимается Docker Engine
        (демон и CLI) и Docker Compose.
  \item[containerd] Среда выполнения контейнеров, которую Docker использует
        внутри. Ею пользуются и Kubernetes, и другие системы.
  \item[runc] Низкоуровневый инструмент, который по стандарту OCI создаёт
        контейнер средствами ядра (namespaces, cgroups).
  \item[Podman] Совместимая с Docker альтернатива без постоянного демона,
        с упором на rootless-режим.
  \item[Kubernetes] Оркестратор: запускает контейнеры на множестве машин,
        перезапускает их, распределяет нагрузку, обновляет без простоя.
        Docker как инструмент сборки образов здесь вполне уместен, хотя
        средой выполнения обычно служит containerd.
\end{description}

\begin{tipbox}
Не пытайтесь выучить Kubernetes, не разобравшись в контейнерах. Ошибки
в Dockerfile, неверные тома и отсутствие проверок здоровья проявляются
в кластере гораздо дороже, чем на локальной машине.
\end{tipbox}

\begin{summarybox}
\begin{itemize}
  \item Контейнер~--- изолированный процесс на общем ядре; образ~---
        неизменяемый шаблон для его запуска.
  \item Виртуальные машины дают сильную изоляцию, но тяжелее и медленнее;
        контейнеры плотнее и быстрее, но делят ядро.
  \item Главная ценность~--- воспроизводимость: один артефакт проходит
        через все среды.
  \item Контейнеры не отменяют резервное копирование и не заменяют
        виртуальные машины там, где нужна жёсткая граница безопасности.
\end{itemize}
\end{summarybox}

\section*{Вопросы для самопроверки}
\begin{selfcheck}
  \item Чем контейнер отличается от виртуальной машины по используемому ядру?
  \item Почему образ называют неизменяемым, а контейнер~--- эфемерным?
  \item Назовите три сценария, где контейнеры дают выигрыш, и два, где
        они могут оказаться неудобны.
  \item Что произойдёт с файлами, созданными внутри контейнера, если
        удалить контейнер без тома?
  \item Зачем нужен стандарт OCI?
\end{selfcheck}
